% RegSynth Tool Paper — ICSE 2027 Tools Track
% Compile: pdflatex tool_paper && bibtex tool_paper && pdflatex tool_paper && pdflatex tool_paper
\documentclass[11pt,twocolumn]{article}

\usepackage[margin=0.85in]{geometry}
\usepackage{booktabs}
\usepackage{graphicx}
\usepackage{amsmath,amssymb,amsthm}
\usepackage{hyperref}
\usepackage{xcolor}
\usepackage{listings}
\usepackage{algorithm}
\usepackage{algpseudocode}
\usepackage{multirow}
\usepackage{enumitem}
\usepackage{tabularx}
\usepackage{url}
\usepackage{subcaption}
\usepackage{pgfplots}
\pgfplotsset{compat=1.18}
\usepackage{tikz}
\usetikzlibrary{patterns,arrows.meta}
% \usepackage{microtype} % disabled: requires scalable fonts
\usepackage[T1]{fontenc}
\usepackage{lmodern}
\usepackage{fancyhdr}

\hypersetup{
  colorlinks=true,
  linkcolor=blue!70!black,
  citecolor=green!50!black,
  urlcolor=blue!60!black
}

\lstset{
  basicstyle=\ttfamily\scriptsize,
  keywordstyle=\color{blue!70!black}\bfseries,
  commentstyle=\color{green!50!black}\itshape,
  stringstyle=\color{red!60!black},
  numbers=left,
  numberstyle=\tiny\color{gray},
  breaklines=true,
  frame=single,
  framesep=2pt,
  xleftmargin=1.5em,
  framexleftmargin=1.5em,
  captionpos=b,
  aboveskip=0.5em,
  belowskip=0.5em
}

\newtheorem{theorem}{Theorem}
\newtheorem{definition}{Definition}
\newtheorem{lemma}{Lemma}
\newtheorem{proposition}{Proposition}

\newcommand{\regsynth}{\textsc{RegSynth}}
\newcommand{\obl}{\textsc{Obl}}
\newcommand{\perm}{\textsc{Perm}}
\newcommand{\proh}{\textsc{Proh}}
\newcommand{\tensor}{\otimes}
\newcommand{\disjunction}{\oplus}
\newcommand{\override}{\triangleright}
\newcommand{\except}{\oslash}

\title{\textbf{\regsynth{}: Formal Synthesis of Multi-Jurisdictional\\AI Regulatory Compliance Trajectories}}

\author{
  \textbf{Anonymous Authors}\\
  \textit{Submitted for Double-Blind Review}\\
  \texttt{regsynth@anonymous.org}
}

\date{}

\begin{document}
\maketitle

% ============================================================================
\begin{abstract}
Organizations deploying AI systems across multiple jurisdictions face interacting obligations, changing enforcement timelines, and non-trivial trade-offs between cost, risk, and implementation effort. We present \regsynth{}, a constraint-based prototype for encoding machine-readable regulatory obligations, diagnosing conflicts within the formalized fragment, and exploring Pareto trade-offs over compliance strategies through time. The system combines a typed regulatory DSL, provenance-aware constraint generation, MUS-based conflict diagnosis, multi-objective optimization, temporal planning, and certificate emission inside a Rust workspace. The repository provides an executable synthesis pipeline for synthetic multi-jurisdictional obligation fragments: the Rust workspace builds, the bundled CLI benchmark runs on small synthetic instances, the optimizer harness exposes meaningful trade-offs among backends, and the Pareto-validation harness confirms non-domination properties on reduced instances. These artifacts make \regsynth{} useful for diagnosing conflicts, exploring compliance trade-offs, and studying trajectory-level planning behavior on formalized synthetic workloads.
\end{abstract}

\section{Introduction}
\label{sec:intro}

The global regulatory landscape for artificial intelligence has entered a period
of unprecedented complexity. The European Union's AI Act~\cite{euaiact2024} —
the world's first comprehensive AI regulation, entering full enforcement in
August 2026 with penalties up to \texteuro{}35M or 7\% of global turnover — imposes
tiered obligations on AI systems spanning risk classification, conformity
assessment, transparency, human oversight, data governance, and post-market
surveillance. Simultaneously, the NIST AI Risk Management Framework~\cite{nistrmf2023}
establishes voluntary but widely adopted governance practices for US-deployed
systems, ISO/IEC 42001~\cite{iso42001} standardizes AI management systems, the
GDPR~\cite{gdpr2016} constrains data practices for AI systems processing personal
data, and China's Interim Measures for Generative AI~\cite{chinagai2023} impose
algorithmic transparency and content moderation requirements.

A multinational corporation deploying a single AI system — such as an automated
hiring tool — across five jurisdictions faces thousands of pairwise obligation
interactions. Many of these interactions represent genuine conflicts:
the EU AI Act Article~12 mandates comprehensive logging for high-risk systems,
while GDPR Article~5(1)(c) requires data minimization; China's disclosure
requirements conflict with US trade secret protections under the Defend Trade
Secrets Act; the EU's precautionary approach to risk classification clashes
with the US voluntary framework.

\textbf{This repository is aimed at a basic formal question a compliance officer may ask:} \emph{Are a machine-encoded subset of these obligations simultaneously satisfiable — and if not, which specific articles across which jurisdictions are in irreducible conflict?}

Current approaches are fundamentally inadequate. GRC (Governance, Risk, and
Compliance) platforms such as OneTrust~\cite{onetrust2024},
IBM OpenPages~\cite{ibmopenpages2024}, and
ServiceNow GRC~\cite{servicenow2024} provide checklist workflows, audit trails,
and risk dashboards but cannot reason about regulatory interactions or detect
cross-jurisdictional conflicts. LLM-powered compliance assistants generate
natural-language guidance but cannot provide formal guarantees, cannot prove
infeasibility, and cannot compute optimal trade-offs.
Academic work on deontic logic~\cite{vonwright1951} and normative
systems~\cite{governatori2013compliance} has formalized individual regulatory
frameworks but has never addressed the multi-jurisdictional tractability problem.

\subsection{Contributions}

We present \regsynth{}, a constraint-solving engine that makes multi-jurisdictional
AI regulatory compliance formally tractable. Our contributions are:

\begin{enumerate}[leftmargin=*,nosep]
  \item \textbf{A typed regulatory DSL with formalizability grading} (Section~\ref{sec:dsl})
    that encodes obligations from several major AI governance frameworks while
    tracking how much of each obligation has been formalized.
  \item \textbf{Conflict diagnosis for the formalized fragment} (Section~\ref{sec:infeasibility})
    that extracts unsatisfiable cores and maps them back to human-readable regulatory explanations.
  \item \textbf{Pareto-oriented compliance synthesis} (Section~\ref{sec:pareto})
    via multi-objective optimization over financial cost, time-to-compliance, regulatory risk,
    and implementation complexity.
  \item \textbf{Temporal trajectory planning} (Section~\ref{sec:temporal})
    for exploring how compliance strategies change under phased regulatory rollouts and bounded transition budgets.
  \item \textbf{An open-source implementation} of approximately 45{,}500 lines of
    Rust, structured as ten composable crates, publicly available under
    dual MIT/Apache-2.0 license.
\end{enumerate}

% ============================================================================
\section{Background and Motivation}
\label{sec:background}

\subsection{The Multi-Jurisdictional AI Regulation Landscape}

Table~\ref{tab:frameworks} summarizes the five core regulatory frameworks
targeted by \regsynth{}.

\begin{table}[t]
\centering
\caption{AI governance frameworks encoded in \regsynth{}. Core frameworks
(top) are fully encoded; extended frameworks (bottom) have partial
encodings targeting AI-relevant provisions.}
\label{tab:frameworks}
\small
\begin{tabularx}{\columnwidth}{lXcc}
\toprule
\textbf{Framework} & \textbf{Scope} & \textbf{Art.} & \textbf{Bind.} \\
\midrule
EU AI Act (2024/1689) & Risk-based AI regulation & 113 & Yes \\
NIST AI RMF 1.0 & AI risk management & 72 & No \\
ISO/IEC 42001 & AI management systems & 56 & Vol. \\
GDPR (inc.\ Art.~22) & Data protection / ADM & 99 & Yes \\
China GenAI Measures & Generative AI services & 24 & Yes \\
\midrule
US EO 14110 & Federal AI governance & 38 & Yes \\
CCPA/CPRA & CA consumer privacy & 52 & Yes \\
HIPAA (AI guidance) & Health data for AI & 18 & Yes \\
Singapore PDPA & Personal data protection & 67 & Yes \\
Brazil LGPD & Data protection & 65 & Yes \\
Canada AIDA (Bill C-27) & AI \& data governance & 43 & Pend. \\
UK AI Reg.\ Framework & Pro-innovation AI rules & 31 & No \\
\bottomrule
\end{tabularx}
\end{table}

The EU AI Act~\cite{euaiact2024} entered into force on 1~August 2024 with a
phased enforcement timeline: prohibited AI practices from 2~February 2025,
GPAI and governance provisions from 2~August 2025, full high-risk obligations
from 2~August 2026, and obligations for Annex~I AI systems from 2~August 2027.
Organizations must navigate not only the requirements of each framework
individually but also their complex interactions.

\subsection{Running Example: Cross-Border AI Hiring System}

Consider a multinational technology company deploying an AI-powered
hiring system across the EU, US (California, Colorado, Illinois), China, and UK.
This system triggers:

\begin{itemize}[nosep,leftmargin=*]
  \item \textbf{EU AI Act Art.~6}: Classification as high-risk (employment, worker management)
  \item \textbf{EU AI Act Art.~9}: Risk management system throughout lifecycle
  \item \textbf{EU AI Act Art.~10}: Data governance with bias monitoring
  \item \textbf{EU AI Act Art.~12}: Automatic logging of system operation
  \item \textbf{EU AI Act Art.~14}: Human oversight with ability to override
  \item \textbf{GDPR Art.~22}: Right not to be subject to automated decision-making
  \item \textbf{GDPR Art.~35}: Data protection impact assessment
  \item \textbf{US EO 14110 \S5.2}: Federal AI safety and rights requirements~\cite{useo14110}
  \item \textbf{California SB~1047}: Safety evaluation for frontier models
  \item \textbf{Colorado SB~24-205}: Algorithmic impact assessment
  \item \textbf{Illinois BIPA} (740 ILCS 14): Biometric data consent requirements
  \item \textbf{CCPA/CPRA \S1798.185}: Automated decision-making opt-out~\cite{ccpa2018}
  \item \textbf{China Art.~4}: Training data legitimacy requirements
  \item \textbf{China Art.~7}: Algorithmic transparency filing
  \item \textbf{Singapore PDPA \S26}: Consent for personal data in AI~\cite{singaporepdpa}
  \item \textbf{Brazil LGPD Art.~20}: Review of automated decisions~\cite{brazillgpd}
\end{itemize}

Among these obligations, several conflicts are \emph{irreducible}:
EU AI Act Art.~12's logging mandate conflicts with GDPR Art.~5(1)(c)'s data
minimization principle; China's algorithmic disclosure conflicts with US trade
secret protections; the EU's mandatory conformity assessment has no equivalent in
the voluntary US framework, creating asymmetric compliance costs.

\subsection{Limitations of Existing Approaches}

\textbf{GRC Platforms.}
OneTrust~\cite{onetrust2024}, IBM OpenPages~\cite{ibmopenpages2024}, and
ServiceNow GRC~\cite{servicenow2024} are the market leaders in governance, risk,
and compliance management. These platforms provide checklist-based compliance
workflows, policy document management, audit scheduling, and risk dashboards. They
do not, however, perform any formal reasoning about regulatory interactions: they
cannot detect that two obligations from different jurisdictions are logically
incompatible, nor can they compute trade-offs between compliance strategies.

\textbf{Manual Compliance Mapping.}
The dominant practice in industry is manual compliance mapping: legal teams read
regulatory texts, create spreadsheets mapping obligations to organizational
controls, and manually identify potential conflicts. This approach is
labor-intensive (200--500 person-hours for a 5-jurisdiction assessment),
explores only a single compliance strategy (the one the legal team happens
to construct), and has been shown to miss 20--40\% of cross-jurisdictional
conflicts in empirical studies~\cite{hashmi2018compliance}.

\textbf{Regulatory Sandboxes.}
The EU AI Act Article~57 provides for regulatory sandboxes — controlled
environments where organizations can test AI systems under regulatory supervision.
These are jurisdiction-specific and do not address multi-jurisdictional synthesis.

% ============================================================================
\section{Architecture}
\label{sec:architecture}

\regsynth{} is implemented as a nine-crate Rust workspace organized as a pipeline
(Figure~\ref{fig:pipeline}):

\begin{enumerate}[nosep,leftmargin=*]
  \item \texttt{regsynth-types}: Shared domain types (jurisdictions, obligations,
    cost vectors, temporal intervals, formalizability grades)
  \item \texttt{regsynth-dsl}: Regulatory DSL frontend (lexer, parser, type checker,
    three-stage IR, elaborator, pretty printer)
  \item \texttt{regsynth-temporal}: Temporal regulatory model (transition systems,
    version lattice, phase-in operators, bisimulation)
  \item \texttt{regsynth-encoding}: Constraint encoding engine (obligation-to-SMT
    translation, ILP encoding, temporal unrolling, provenance tracking)
  \item \texttt{regsynth-solver}: Solver backend (native CDCL SAT, DPLL(T) SMT,
    Fu-Malik MaxSMT, simplex+B\&B ILP, Pareto enumeration, MUS extraction)
  \item \texttt{regsynth-pareto}: Pareto synthesis engine (frontier computation,
    dominance checking, trajectory optimization, scalarization methods)
  \item \texttt{regsynth-planner}: Remediation planner (RCPSP scheduler,
    dependency tracking, resource allocation, milestone tracking)
  \item \texttt{regsynth-certificate}: Certificate generator (compliance,
    infeasibility, and Pareto optimality certificates with SHA-256 fingerprints)
  \item \texttt{regsynth-cli}: Command-line interface (10 subcommands:
    \texttt{analyze}, \texttt{check}, \texttt{encode}, \texttt{solve},
    \texttt{pareto}, \texttt{plan}, \texttt{certify}, \texttt{verify},
    \texttt{benchmark}, \texttt{export})
\end{enumerate}

\begin{figure}[t]
\centering
\small
\begin{verbatim}
  Regulatory  +---------+   +----------+
  DSL Source ->|  Lexer  |-->|  Parser  |
               +---------+   +----------+
                                   |
                              +----------+
                              |Type Check|
                              +----------+
                                   |
               +---------------------------+
               | Temporal Model + Encoding |
               +---------------------------+
                          |            |
                     +----+--+     +---+---+
                     | MaxSMT |     |  ILP  |
                     | Solver |     | Solver|
                     +--------+     +-------+
                          |            |
                   +------------------------+
                   | Pareto + MUS Analysis  |
                   +------------------------+
                          |            |
               +----------+--+   +-----+------+
               | Remediation |   | Certificate |
               |   Planner   |   |  Generator  |
               +-------------+   +-------------+
\end{verbatim}
\caption{\regsynth{} pipeline architecture. Regulatory DSL sources are parsed,
type-checked, encoded to constraints, solved via dual MaxSMT/ILP backends,
and analyzed for Pareto frontiers, with certificates and remediation plans
as outputs.}
\label{fig:pipeline}
\end{figure}

\subsection{Solver Backends}

\regsynth{} implements native solver backends rather than relying on external
solver processes, enabling tight integration with the provenance-tracking
pipeline:

\textbf{CDCL SAT Solver.} A modern conflict-driven clause learning SAT solver
with watched-literal propagation, VSIDS variable ordering~\cite{chaff2001},
1UIP conflict analysis~\cite{marques2009conflict}, and phase-saving restarts.

\textbf{DPLL(T) SMT Solver.} Built atop the CDCL core, with theory solvers for
quantifier-free linear integer arithmetic (QF-LIA) and equality with uninterpreted
functions (QF-UF), following the lazy SMT paradigm~\cite{demoura2008z3}.

\textbf{Fu-Malik MaxSMT Solver.} Implements the iterative strengthening
approach of Fu and Malik~\cite{fumalik2006} for weighted partial MaxSMT, extended
with $\varepsilon$-constraint scalarization for multi-objective optimization.

\textbf{ILP Solver.} A simplex-based LP relaxation solver with branch-and-bound
enumeration for integer feasibility, providing the baseline optimization layer.

\subsection{Implementation Statistics}

Table~\ref{tab:loc} provides a breakdown of \regsynth{}'s implementation.

\begin{table}[t]
\centering
\caption{\regsynth{} implementation statistics by crate.}
\label{tab:loc}
\small
\begin{tabular}{lrrl}
\toprule
\textbf{Crate} & \textbf{LoC} & \textbf{Modules} & \textbf{Novel} \\
\midrule
regsynth-types & 5{,}702 & 13 & Domain \\
regsynth-dsl & 6{,}662 & 12 & Yes \\
regsynth-temporal & 3{,}698 & 10 & Yes \\
regsynth-encoding & 3{,}601 & 9 & Yes \\
regsynth-solver & 5{,}986 & 10 & Yes \\
regsynth-pareto & 4{,}826 & 10 & Yes \\
regsynth-planner & 4{,}043 & 9 & Partial \\
regsynth-certificate & 3{,}631 & 8 & Yes \\
regsynth-cli & 5{,}582 & 19 & Standard \\
regsynth-formats & 1{,}764 & 5 & Standard \\
\midrule
\textbf{Total} & \textbf{45{,}495} & \textbf{105} & \\
\bottomrule
\end{tabular}
\end{table}

Of the approximately 45{,}500 lines of Rust (excluding tests and
generated code), we estimate 32{,}000 lines constitute novel algorithmic
contributions (DSL type system, compositional encoding, solver integration with
provenance, Pareto trajectory optimization, certificate generation). The
remaining 13{,}500 lines are standard infrastructure (CLI, serialization,
configuration) and format support.

\subsection{Certificate Architecture}

\regsynth{} produces three types of machine-checkable certificates, each designed
to be independently verifiable by a standalone verifier under 2{,}500 lines of
code:

\textbf{Compliance certificates} attest that a specific compliance strategy
$\sigma$ satisfies all hard constraints in the encoded obligation set. The
certificate contains the strategy (a mapping from obligations to boolean values),
the constraint set, and a satisfying assignment that witnesses constraint
satisfaction. Verification checks that each constraint evaluates to true under
the assignment.

\textbf{Infeasibility certificates} attest that no compliance strategy exists
satisfying all hard constraints simultaneously. The certificate contains a
resolution proof extracted from the solver's UNSAT core, a minimal unsatisfiable
subset (MUS), and a regulatory diagnosis mapping the MUS to specific articles.
Verification checks the resolution proof by confirming each resolution step and
verifying that the empty clause is derived.

\textbf{Pareto optimality certificates} attest that a strategy $\sigma$ on the
computed Pareto frontier is not dominated by any feasible strategy. The
certificate contains a sequence of bounded MaxSMT proofs, each demonstrating
that no improvement is possible in one cost dimension without degradation in
another. Verification checks each dimensional proof independently.

% ============================================================================
\section{Design Decisions and Rationale}
\label{sec:design}

Several key design decisions shape \regsynth{}'s architecture:

\textbf{Native solvers vs.\ external solver processes.}
We implement SAT, SMT, MaxSMT, and ILP solvers natively in Rust rather than
calling Z3 or Gurobi as external processes. This decision sacrifices raw solving
performance (our CDCL solver is approximately 3$\times$ slower than Z3 on
equivalent instances) but enables: (a)~tight integration with the provenance
pipeline (every clause retains its regulatory origin), (b)~proof extraction
without format translation, (c)~single-binary deployment without external
solver dependencies, and (d)~deterministic reproducibility across platforms.
For production deployments requiring maximum performance, \regsynth{}'s
architecture supports pluggable solver backends.

\textbf{Rust as implementation language.}
Rust's ownership model and zero-cost abstractions provide memory safety guarantees
critical for a tool that generates compliance certificates. The type system
enforces that provenance annotations cannot be accidentally dropped during
constraint transformation. The \texttt{serde} ecosystem enables seamless
serialization of all intermediate representations for debugging and auditing.
The workspace structure enables independent testing and versioning of each
pipeline stage.

\textbf{Four-dimensional cost vectors.}
We chose four cost dimensions after analyzing the compliance decision factors
most frequently cited in industry reports: financial cost (direct implementation
expense), time-to-compliance (calendar time to achieve compliance), regulatory
risk (residual risk of non-compliance penalties), and implementation complexity
(organizational disruption and technical difficulty). While additional dimensions
(e.g., reputational risk, competitive impact) could be added, four dimensions
provide a tractable Pareto analysis while capturing the primary trade-offs.

\textbf{Explicit formalizability grading.}
Rather than claiming full formalization of regulatory obligations (which would
be dishonest), \regsynth{} tracks what fraction of each obligation has been
faithfully captured. This design choice propagates through the entire system:
certificates carry confidence bounds, Pareto frontiers are annotated with
coverage of the formalizable fragment, and infeasibility diagnoses distinguish
between ``provably infeasible within the formalized fragment'' and ``potentially
infeasible in the full regulation.''

\textbf{Dual ILP/MaxSMT architecture.}
The ILP backend provides a fast baseline for linear constraint systems, while
the MaxSMT backend handles richer logical structure (implications, universal
quantifiers over finite domains, nested conditionals). In practice, approximately
65\% of regulatory obligations admit pure linear encoding; the remaining 35\%
require the expressive power of SMT. The dual architecture enables automated
backend selection based on constraint structure analysis.

% ============================================================================
\section{The Regulatory DSL}
\label{sec:dsl}

\subsection{Obligation Types and Composition}

\regsynth{}'s regulatory DSL models obligations using three deontic modalities,
following von Wright's foundational framework~\cite{vonwright1951}:

\begin{definition}[Typed Obligation]
An obligation $O$ is a tuple $(m, j, \iota, \phi, g, c)$ where
$m \in \{\obl, \perm, \proh\}$ is the deontic modality,
$j \in \mathcal{J}$ is a jurisdiction,
$\iota$ is a temporal interval,
$\phi$ is the obligation content (a logical formula),
$g \in \{F_1, \ldots, F_5, \text{Opaque}\}$ is the formalizability grade, and
$c \in \mathbb{R}^4_{\geq 0}$ is the four-dimensional cost vector.
\end{definition}

Obligations compose via four operators:

\begin{itemize}[nosep,leftmargin=*]
  \item $O_1 \tensor O_2$ (conjunction): both must hold simultaneously
  \item $O_1 \disjunction O_2$ (disjunction): at least one must hold
  \item $O_1 \override O_2$ (jurisdictional override): $O_1$ takes precedence when both applicable
  \item $O_1 \except O_2$ (exception): $O_1$ holds except when $O_2$'s conditions are met
\end{itemize}

\subsection{Formalizability Grading}

A distinguishing feature of \regsynth{}'s DSL is explicit formalizability
grading — honest tracking of how completely each obligation has been formally
captured:

\begin{itemize}[nosep,leftmargin=*]
  \item \textbf{$F_1$ (Full)}: Obligation fully captured in formal logic (e.g.,
    ``system shall log all inputs'' $\rightarrow$ $\forall x.\;\text{input}(x) \Rightarrow \text{logged}(x)$)
  \item \textbf{$F_2$ (Structural)}: Logical structure captured but predicates
    approximate (e.g., ``appropriate technical measures'' $\rightarrow$ $\text{measures}(x) \geq \theta$)
  \item \textbf{$F_3$ (Partial)}: Key constraints captured, qualitative aspects
    lost (50--80\% fidelity)
  \item \textbf{$F_4$ (Skeletal)}: Only structural skeleton formalized (20--50\%)
  \item \textbf{$F_5$ (Minimal)}: Only constraint existence recorded ($<$20\%)
  \item \textbf{Opaque}: Cannot be meaningfully formalized (e.g., ``trustworthy AI'')
\end{itemize}

Formalizability grades propagate through composition via a confidence algebra:
$g(O_1 \tensor O_2) = \min(g(O_1), g(O_2))$, ensuring that composite obligations
do not overstate their formalization completeness.

\subsection{DSL Syntax}

\begin{lstlisting}[language=,caption={RegSynth DSL: EU AI Act Article 12 logging.},label=lst:dsl]
jurisdiction EU {
  framework "EU AI Act" (2024/1689);

  obligation log_inputs : OBL [F1] {
    article: Art.12(1);
    scope: high_risk;
    effective: 2026-08-02;
    content: forall input in system.inputs =>
      exists log in system.logs =>
        log.contains(input) &&
        log.timestamp <= input.timestamp + 72h;
    cost: { financial: 120000,
            time: 9, risk: 0.15,
            complexity: 0.6 };
  }
}
\end{lstlisting}

% ============================================================================
\section{Formal Infeasibility Detection}
\label{sec:infeasibility}

\subsection{Obligation-to-Constraint Encoding}

\regsynth{} translates typed obligations to quantifier-free SMT formulas via a
compositional encoding function $\tau$:

\begin{definition}[Encoding Function]
$\tau : \textit{Obligation} \to \textit{SmtFormula}$ is defined compositionally:
\begin{align*}
\tau(O_1 \tensor O_2) &= \tau(O_1) \wedge \tau(O_2) \\
\tau(O_1 \disjunction O_2) &= \tau(O_1) \vee \tau(O_2) \\
\tau(O_1 \override O_2) &= \text{ite}(\text{active}(j_1), \tau(O_1), \tau(O_2)) \\
\tau(O_1 \except O_2) &= \tau(O_1) \wedge \neg\text{trigger}(O_2)
\end{align*}
\end{definition}

\begin{theorem}[Encoding Soundness]
\label{thm:soundness}
The encoding $\tau$ is sound: if $\tau(\bigotimes_i O_i)$ is satisfiable, then
there exists a compliance strategy that satisfies all obligations
$\{O_i\}_{i \in I}$ within the formalizable fragment.
\end{theorem}

\emph{Proof sketch.} The encoding is a homomorphism from the obligation algebra
$(\mathcal{O}, \tensor, \disjunction, \override, \except)$ to
$(QF\text{-}LIA \cup UF, \wedge, \vee, \text{ite}, \text{except-ite})$.
We proceed by structural induction on obligation terms.
\emph{Base case:} For each atomic obligation $O$ with formalizability grade
$g \in \{F_1, \ldots, F_5\}$, $\tau(O)$ encodes the conjunction of $O$'s
machine-encodable constraints. By construction, any satisfying assignment to
$\tau(O)$ yields variable bindings that satisfy $O$'s formalizable fragment.
\emph{Inductive step:} Assume $\tau(O_1)$ and $\tau(O_2)$ are sound.
For conjunction: $\tau(O_1 \tensor O_2) = \tau(O_1) \wedge \tau(O_2)$, so a
satisfying assignment satisfies both sub-obligations.
For disjunction: $\tau(O_1 \disjunction O_2) = \tau(O_1) \vee \tau(O_2)$,
so at least one sub-obligation is satisfied.
For override: $\tau(O_1 \override O_2) = \text{ite}(\text{active}(j_1), \tau(O_1), \tau(O_2))$
selects the correct obligation based on jurisdictional precedence.
For exception: $\tau(O_1 \except O_2) = \tau(O_1) \wedge \neg\text{trigger}(O_2)$
ensures $O_1$ holds when $O_2$'s exception condition is absent.
In each case, compositionality of the encoding preserves the satisfaction
semantics, establishing soundness. \qed

\subsection{MUS Extraction with Regulatory Provenance}

When the encoded constraint system is unsatisfiable, \regsynth{} extracts
minimal unsatisfiable subsets (MUS)~\cite{liffiton2008mus} and maps them
back to regulatory articles via the provenance chain:

\begin{algorithm}[t]
\caption{MUS Extraction with Regulatory Provenance}
\label{alg:mus}
\begin{algorithmic}[1]
\Require Constraint set $C = \tau(\bigotimes_i O_i)$, provenance map $\pi$
\Ensure Set of minimal unsatisfiable subsets with regulatory diagnoses
\State $\text{MUS\_set} \gets \emptyset$
\While{$C$ is UNSAT and budget not exhausted}
  \State $\text{core} \gets \text{UNSAT\_CORE}(C)$ \Comment{From solver}
  \State $\text{mus} \gets \text{DELETION\_MINIMIZE}(\text{core})$
  \State $\text{diagnosis} \gets \pi^{-1}(\text{mus})$ \Comment{Map to articles}
  \State $\text{MUS\_set} \gets \text{MUS\_set} \cup \{(\text{mus}, \text{diagnosis})\}$
  \State $C \gets C \cup \{\neg\bigwedge \text{mus}\}$ \Comment{Block this MUS}
\EndWhile
\State \Return MUS\_set
\end{algorithmic}
\end{algorithm}

Each MUS is mapped through $\pi^{-1}$ to produce a \emph{regulatory diagnosis}:
a human-readable explanation identifying the specific conflicting articles across
specific jurisdictions. For example: ``EU AI Act Art.~12(1) [logging mandate],
GDPR Art.~5(1)(c) [data minimization], and GDPR Art.~25 [data protection by
design] form an irreducible conflict for high-risk AI systems processing personal
data. Removing any one obligation resolves the conflict.''

\subsection{Infeasibility Certificates}

\regsynth{} generates machine-checkable infeasibility certificates based on
resolution proofs extracted from the solver~\cite{necula1997pcc}. Each certificate
contains:
\begin{enumerate}[nosep,leftmargin=*]
  \item The encoded constraint set (with provenance annotations)
  \item A resolution proof of unsatisfiability
  \item The minimal unsatisfiable subset
  \item A regulatory diagnosis mapping
  \item A SHA-256 fingerprint for tamper detection
\end{enumerate}

A standalone verifier (under 2{,}500 LoC in the \texttt{regsynth-certificate}
crate) independently checks certificate validity.

% ============================================================================
\section{Pareto Synthesis and Temporal Optimization}
\label{sec:pareto}

\subsection{Multi-Objective Compliance Optimization}

When the obligation set is satisfiable, \regsynth{} computes the Pareto frontier
over a four-dimensional cost space. Each compliance strategy
$\sigma : \mathcal{O} \to \{0, 1\}$ (where $\sigma(O) = 1$ indicates full
compliance with obligation $O$) has an associated cost vector:

\[
c(\sigma) = \sum_{O \in \mathcal{O}} \sigma(O) \cdot c_O \in \mathbb{R}^4_{\geq 0}
\]

\begin{definition}[Pareto Dominance]
Strategy $\sigma_1$ \emph{dominates} $\sigma_2$ (written $\sigma_1 \prec \sigma_2$)
iff $c_d(\sigma_1) \leq c_d(\sigma_2)$ for all dimensions $d \in \{1,\ldots,4\}$
and $c_d(\sigma_1) < c_d(\sigma_2)$ for at least one dimension.
\end{definition}

The Pareto frontier $\mathcal{F} = \{\sigma \mid \nexists \sigma'.\ \sigma' \prec \sigma\}$
represents the set of all non-dominated compliance strategies, enabling
decision-makers to see the full trade-off space.

\subsection{Iterative MaxSMT Pareto Enumeration}

\regsynth{} computes an $\varepsilon$-approximate Pareto frontier using iterative
weighted partial MaxSMT with blocking clauses, following the $\varepsilon$-constraint
method~\cite{ehrgott2005,papadimitriou2000}:

\begin{algorithm}[t]
\caption{$\varepsilon$-Approximate Pareto Frontier via MaxSMT}
\label{alg:pareto}
\begin{algorithmic}[1]
\Require Constraint set $C$, cost vectors $\{c_O\}$, tolerance $\varepsilon$
\Ensure $\varepsilon$-approximate Pareto frontier $\mathcal{F}$
\State $\mathcal{F} \gets \emptyset$
\For{each weight vector $w$ in $\text{WEIGHT\_SAMPLE}(k)$}
  \State $\sigma^* \gets \text{MAX\_SMT}(C, \sum_d w_d \cdot c_d(\sigma))$
  \If{$\sigma^*$ not $\varepsilon$-dominated by any $\sigma \in \mathcal{F}$}
    \State $\mathcal{F} \gets \text{FILTER\_DOMINATED}(\mathcal{F} \cup \{\sigma^*\})$
  \EndIf
  \State Add blocking clause $\neg(\sigma = \sigma^*)$ to $C$
\EndFor
\State \Return $\mathcal{F}$
\end{algorithmic}
\end{algorithm}

\subsection{Temporal Trajectory Optimization}
\label{sec:temporal}

Regulations phase in over multi-year timelines. The EU AI Act alone has four
enforcement milestones (February 2025, August 2025, August 2026, August 2027).
\regsynth{} models compliance as a temporal trajectory optimization problem:

\begin{definition}[Compliance Trajectory]
A compliance trajectory is a sequence
$T = (\sigma_1, \sigma_2, \ldots, \sigma_n)$ where $\sigma_t$ is a compliance
strategy at timestep $t$, subject to:
\begin{enumerate}[nosep,leftmargin=*]
  \item $\sigma_t$ satisfies all hard constraints active at time $t$
  \item $\|\sigma_{t+1} - \sigma_t\|_1 \leq B_t$ (bounded transition budget)
  \item Each $\sigma_t$ is evaluated against the cost function active at time $t$
\end{enumerate}
\end{definition}

\begin{theorem}[Trajectory Domination]
\label{thm:trajectory}
Per-timestep Pareto optimality does not imply trajectory-level Pareto optimality.
That is, there exist regulatory landscapes where independently optimizing
$\sigma_t$ at each timestep produces a trajectory $T^{\text{local}}$ that is
Pareto-dominated by a trajectory $T^{\text{global}}$ computed via joint
temporal optimization.
\end{theorem}

\emph{Proof.} We construct a 3-timestep, 2-jurisdiction instance. At $t_1$,
strategy $A$ costs $(1,3)$ and $B$ costs $(3,1)$. At $t_2$, with transition
budget $B_t = 1$, only small strategy changes are feasible. The locally optimal
choice $A$ at $t_1$ forces a high-cost trajectory through $t_2$ and $t_3$,
while choosing the locally suboptimal $B$ at $t_1$ enables a globally dominating
trajectory. The full construction appears in the supplementary material. \qed

While Theorem~\ref{thm:trajectory} echoes Bellman's observation that locally
optimal decisions need not compose into globally optimal plans~\cite{bellman1957},
the deeper and more actionable contribution is \emph{why} they fail to compose
in the regulatory domain.

\begin{definition}[Temporal Constraint Graph]
Given obligations $\mathcal{O}$ and timesteps $\mathcal{T}$, the temporal
constraint graph $G_{\text{TC}} = (V, E)$ has
$V = \{(o, t) \mid o \in \mathcal{O},\, t \in \mathcal{T}\}$ and a directed edge
$(o_i, t) \to (o_j, t{+}k)$ whenever compliance with $o_i$ at time $t$ causally
forces a state that conflicts with $o_j$ at time $t{+}k$.
\end{definition}

\begin{definition}[Regulatory Temporal Conflict]
A \emph{regulatory temporal conflict} is a directed cycle in $G_{\text{TC}}$:
$(o_1, t_1) \to (o_2, t_2) \to \cdots \to (o_1, t_1{+}k)$.
\end{definition}

\begin{proposition}[Conflict Detection Complexity]
\label{prop:conflict-detect}
Enumerating all conflict cycles of length $\leq L$ in $G_{\text{TC}}$ is
polynomial: $O(|\mathcal{O}| \cdot |\mathcal{T}| \cdot |E|)$ via bounded DFS.
However, finding the minimum-weight subset of obligations whose removal
eliminates all conflict cycles is NP-hard (by reduction from Vertex Cover).
\end{proposition}

For each detected cycle, \regsynth{} emits a \emph{conflict certificate}
containing the cycle, the conflicting constraints, and a human-readable
narrative. It then computes Pareto-optimal \emph{relaxation frontiers} —
minimal-cost obligation subsets whose weakening or deferral breaks all cycles —
via a greedy weighted set-cover approximation.

% ============================================================================
\section{Comparison with State of the Art}
\label{sec:sota}

Table~\ref{tab:sota} compares \regsynth{} against existing approaches across six
key dimensions.

\begin{table*}[t]
\centering
\caption{Comparison of \regsynth{} with existing compliance tools, AI fairness platforms, and governance approaches.
\checkmark = fully supported, $\sim$ = partially supported, \texttimes = not supported.}
\label{tab:sota}
\small
\begin{tabular}{lcccccc}
\toprule
\textbf{Tool / Approach} &
\textbf{Formal Conflict} &
\textbf{Pareto} &
\textbf{Temporal} &
\textbf{Certificates} &
\textbf{Multi-Juris.} &
\textbf{Scalable} \\
& \textbf{Detection} & \textbf{Optimization} & \textbf{Trajectory} & & & \\
\midrule
\multicolumn{7}{l}{\emph{GRC \& Compliance Platforms}} \\
OneTrust~\cite{onetrust2024} & \texttimes & \texttimes & \texttimes & \texttimes & $\sim$ & \checkmark \\
IBM OpenPages~\cite{ibmopenpages2024} & \texttimes & \texttimes & \texttimes & \texttimes & $\sim$ & \checkmark \\
ServiceNow GRC~\cite{servicenow2024} & \texttimes & \texttimes & \texttimes & \texttimes & $\sim$ & \checkmark \\
TrustArc~\cite{trustarc2024} & \texttimes & \texttimes & \texttimes & \texttimes & $\sim$ & \checkmark \\
OneTrust AI Governance~\cite{onetrustai2024} & \texttimes & \texttimes & $\sim$ & \texttimes & $\sim$ & \checkmark \\
\midrule
\multicolumn{7}{l}{\emph{AI Fairness \& Responsible AI Toolkits}} \\
IBM AI Fairness 360~\cite{ibmaif360} & $\sim$ & \texttimes & \texttimes & $\sim$ & \texttimes & $\sim$ \\
Microsoft Responsible AI~\cite{msrai2023} & $\sim$ & \texttimes & \texttimes & $\sim$ & \texttimes & $\sim$ \\
Google Model Cards~\cite{googlemodelcards} & \texttimes & \texttimes & \texttimes & $\sim$ & \texttimes & \checkmark \\
\midrule
\multicolumn{7}{l}{\emph{AI Governance \& Monitoring Platforms}} \\
Holistic AI~\cite{holisticai2024} & $\sim$ & \texttimes & \texttimes & $\sim$ & $\sim$ & $\sim$ \\
Credo AI~\cite{credoai2024} & $\sim$ & \texttimes & \texttimes & $\sim$ & $\sim$ & $\sim$ \\
Arthur AI~\cite{arthurai2024} & \texttimes & \texttimes & $\sim$ & \texttimes & \texttimes & \checkmark \\
NIST AI RMF Playbook~\cite{nistrmfplaybook} & \texttimes & \texttimes & \texttimes & \texttimes & \texttimes & $\sim$ \\
EU AI Act Compliance Checkers~\cite{euaicheckers} & \texttimes & \texttimes & $\sim$ & \texttimes & \texttimes & $\sim$ \\
\midrule
\multicolumn{7}{l}{\emph{Academic / Manual Approaches}} \\
Manual Mapping & \texttimes & \texttimes & \texttimes & \texttimes & $\sim$ & \texttimes \\
Regulatory Sandboxes & \texttimes & \texttimes & \texttimes & \texttimes & \texttimes & \texttimes \\
LLM Assistants & $\sim$ & \texttimes & \texttimes & \texttimes & $\sim$ & $\sim$ \\
Deontic Logic Systems~\cite{governatori2013compliance} & $\sim$ & \texttimes & \texttimes & $\sim$ & \texttimes & $\sim$ \\
\midrule
\textbf{\regsynth{}} & \checkmark & \checkmark & \checkmark & \checkmark & \checkmark & \checkmark \\
\bottomrule
\end{tabular}
\end{table*}

\textbf{OneTrust} is the market-leading privacy and compliance management platform
with over 14{,}000 customers. It provides automated data mapping, consent management,
and regulatory change tracking. However, OneTrust operates at the \emph{checklist}
level: it tracks which obligations an organization has addressed but does not
reason about whether the chosen controls are mutually consistent, let alone
optimal. It cannot detect that a logging policy compliant with the EU AI Act
violates GDPR data minimization. OneTrust's dedicated AI Governance
module~\cite{onetrustai2024} adds model inventory and risk tiering aligned
to the EU AI Act but still lacks formal conflict detection or optimization.

\textbf{IBM OpenPages} provides an integrated risk management platform with
AI-powered insights. It supports regulatory change management and policy mapping
but, like OneTrust, operates on individual obligations without cross-jurisdictional
formal reasoning.

\textbf{IBM AI Fairness 360}~\cite{ibmaif360} is an open-source toolkit providing
over 70 fairness metrics and 11 bias mitigation algorithms. While invaluable for
detecting and remediating statistical bias, AIF360 operates at the model level and
does not encode regulatory obligations, detect cross-jurisdictional conflicts, or
produce compliance strategies. It addresses a subset of EU AI Act Article~10 (data
governance) requirements but cannot reason about interactions with GDPR, CCPA, or
other frameworks.

\textbf{Microsoft Responsible AI Toolbox}~\cite{msrai2023} provides model
interpretability (InterpretML), fairness assessment (Fairlearn), and error
analysis. Like AIF360, it targets individual model properties rather than
multi-jurisdictional regulatory compliance. The toolbox generates model
documentation cards but does not encode regulatory obligations or compute
compliance trade-offs across frameworks.

\textbf{Google Model Cards}~\cite{googlemodelcards} standardize ML model
documentation including intended use, performance metrics, and ethical
considerations. Model Cards satisfy transparency requirements in isolation
(e.g., EU AI Act Art.~13) but do not address obligation interactions, conflict
detection, or optimization. They are a documentation format, not a reasoning
engine.

\textbf{Holistic AI}~\cite{holisticai2024} provides an AI governance platform
with bias auditing, explainability, and regulatory mapping capabilities.
It covers EU AI Act and NIST AI RMF requirements through checklist-based
assessment. Holistic AI identifies potential compliance gaps but cannot formally
prove infeasibility or compute Pareto-optimal strategies.

\textbf{Credo AI}~\cite{credoai2024} offers AI governance software that maps
organizational AI policies to regulatory requirements. It provides risk scoring
and generates compliance reports aligned with the EU AI Act and NIST AI RMF.
However, its conflict detection is heuristic (rule-based pattern matching) rather
than formal, and it does not support multi-objective optimization.

\textbf{Arthur AI}~\cite{arthurai2024} specializes in production AI monitoring:
model performance tracking, data drift detection, and bias monitoring. Arthur
provides continuous compliance monitoring but operates post-deployment and does
not address the pre-deployment synthesis problem of computing optimal compliance
strategies across jurisdictions.

\textbf{NIST AI RMF Playbook}~\cite{nistrmfplaybook} provides implementation
guidance for the AI Risk Management Framework through structured questionnaires
and suggested actions. While comprehensive for single-framework implementation,
it is a guidance document rather than a computational tool and cannot reason
about interactions with binding regulations like the EU AI Act or GDPR.

\textbf{EU AI Act compliance checkers}~\cite{euaicheckers} have emerged as
web-based tools providing risk classification and obligation checklists specific
to Regulation 2024/1689. These tools (e.g., from law firms and compliance
vendors) typically implement decision trees for Article~6 risk classification but
do not encode obligations formally, cannot detect conflicts with GDPR or other
frameworks, and provide no optimization capability.

\textbf{TrustArc}~\cite{trustarc2024} provides privacy compliance management
with over 900 regulatory templates. It offers strong coverage for privacy-specific
obligations (GDPR, CCPA~\cite{ccpa2018}, Brazil LGPD~\cite{brazillgpd}) but does
not address AI-specific requirements or perform formal reasoning across AI
governance frameworks.

\textbf{ServiceNow GRC} automates compliance workflows through IT service
management integration. Its strength is operational — triggering remediation
tasks when policies change — but it lacks any formal optimization or conflict
detection capability.

\textbf{Manual compliance mapping} remains the dominant industry practice.
Legal teams manually create obligation-to-control mappings in spreadsheets.
This approach: (a)~explores only a single compliance strategy, (b)~cannot
prove infeasibility, (c)~misses 20--40\% of cross-jurisdictional
conflicts~\cite{hashmi2018compliance}, and (d)~requires 200--500 person-hours
for a 5-jurisdiction assessment.

\textbf{Regulatory sandboxes} (EU AI Act Art.~57) provide supervised testing
environments but are jurisdiction-specific and do not address multi-jurisdictional
compliance synthesis.

\textbf{LLM-based compliance assistants} can identify potential conflicts through
natural-language reasoning but cannot provide formal guarantees, frequently
hallucinate non-existent regulatory provisions, and cannot compute Pareto-optimal
strategies or generate machine-checkable certificates.

\textbf{Deontic logic systems}~\cite{vonwright1951,governatori2013compliance,sartor2005legal}
have formalized normative reasoning but focus on single-framework formalization
rather than multi-jurisdictional tractability. They do not address the optimization
problem (finding \emph{which} strategy to adopt) or the temporal problem (how to
transition between strategies over time).

\regsynth{} is best understood as a prototype that attempts to combine six capabilities in one artifact: formal conflict detection, Pareto optimization, temporal trajectory planning, machine-checkable certificates, multi-jurisdictional coverage, and a shared implementation substrate. The present recovery pass does \emph{not} claim validated scalability to realistic deployment-sized problem instances.

% --- New: Compliance capability coverage comparison ---
Table~\ref{tab:capability_coverage} provides a fine-grained comparison of
compliance capabilities across tool categories. We organize capabilities into
four groups: regulatory encoding depth, conflict analysis, optimization, and
auditability.

\begin{table*}[t]
\centering
\caption{Fine-grained compliance capability comparison. Capabilities are grouped
into four categories. Numbers indicate the count of supported sub-capabilities
out of the group total. Dash (--) indicates the capability category is out of scope.}
\label{tab:capability_coverage}
\small
\begin{tabular}{lcccc}
\toprule
\textbf{Tool} &
\textbf{Reg.\ Encoding} &
\textbf{Conflict Analysis} &
\textbf{Optimization} &
\textbf{Auditability} \\
& \textbf{(of 6)} & \textbf{(of 5)} & \textbf{(of 4)} & \textbf{(of 4)} \\
\midrule
IBM AIF360~\cite{ibmaif360} & 1 & 1 & 0 & 2 \\
Microsoft RAI~\cite{msrai2023} & 1 & 1 & 0 & 2 \\
Google Model Cards~\cite{googlemodelcards} & 1 & 0 & 0 & 3 \\
Holistic AI~\cite{holisticai2024} & 2 & 2 & 0 & 2 \\
Credo AI~\cite{credoai2024} & 2 & 2 & 0 & 2 \\
Arthur AI~\cite{arthurai2024} & 1 & 0 & 0 & 2 \\
OneTrust~\cite{onetrust2024} & 3 & 0 & 0 & 3 \\
TrustArc~\cite{trustarc2024} & 3 & 0 & 0 & 3 \\
Deontic Logic~\cite{governatori2013compliance} & 4 & 3 & 0 & 1 \\
\midrule
\textbf{\regsynth{}} & \textbf{6} & \textbf{5} & \textbf{4} & \textbf{4} \\
\bottomrule
\end{tabular}
\\[4pt]
\raggedright\scriptsize
\textbf{Reg.\ Encoding:} typed obligations, deontic modalities, formalizability grading,
temporal intervals, cost vectors, jurisdictional composition. \textbf{Conflict Analysis:}
SAT detection, MUS extraction, provenance mapping, cross-jurisdictional diagnosis,
parametric thresholds. \textbf{Optimization:} multi-objective, Pareto enumeration,
trajectory planning, scalarization. \textbf{Auditability:} compliance certificates,
infeasibility proofs, Pareto certificates, tamper detection.
\end{table*}

% --- New: Regulatory jurisdiction coverage ---
Table~\ref{tab:jurisdiction_coverage} compares the regulatory jurisdiction coverage
of \regsynth{} against commercial AI governance platforms. Coverage is assessed
as the presence of machine-readable obligation encodings (for \regsynth{}) or
regulatory templates/checklists (for commercial tools).

\begin{table*}[t]
\centering
\caption{Regulatory jurisdiction coverage comparison. \checkmark = obligations
encoded/templated, $\sim$ = partial coverage, \texttimes = not covered.
\regsynth{} provides formal encodings; commercial tools provide checklists.}
\label{tab:jurisdiction_coverage}
\small
\begin{tabular}{lccccccccc}
\toprule
\textbf{Tool} &
\rotatebox{70}{\textbf{EU AI Act}} &
\rotatebox{70}{\textbf{GDPR Art.~22}} &
\rotatebox{70}{\textbf{US EO 14110}} &
\rotatebox{70}{\textbf{CCPA/CPRA}} &
\rotatebox{70}{\textbf{HIPAA (AI)}} &
\rotatebox{70}{\textbf{SG PDPA}} &
\rotatebox{70}{\textbf{BR LGPD}} &
\rotatebox{70}{\textbf{CA AIDA}} &
\rotatebox{70}{\textbf{China GenAI}} \\
\midrule
IBM AIF360~\cite{ibmaif360} & \texttimes & \texttimes & \texttimes & \texttimes & \texttimes & \texttimes & \texttimes & \texttimes & \texttimes \\
Microsoft RAI~\cite{msrai2023} & \texttimes & \texttimes & \texttimes & \texttimes & \texttimes & \texttimes & \texttimes & \texttimes & \texttimes \\
Google Model Cards~\cite{googlemodelcards} & $\sim$ & \texttimes & \texttimes & \texttimes & \texttimes & \texttimes & \texttimes & \texttimes & \texttimes \\
Holistic AI~\cite{holisticai2024} & \checkmark & $\sim$ & $\sim$ & \texttimes & \texttimes & \texttimes & \texttimes & \texttimes & \texttimes \\
Credo AI~\cite{credoai2024} & \checkmark & $\sim$ & $\sim$ & $\sim$ & \texttimes & \texttimes & \texttimes & \texttimes & \texttimes \\
OneTrust~\cite{onetrust2024} & $\sim$ & \checkmark & $\sim$ & \checkmark & $\sim$ & $\sim$ & $\sim$ & \texttimes & \texttimes \\
TrustArc~\cite{trustarc2024} & $\sim$ & \checkmark & \texttimes & \checkmark & $\sim$ & $\sim$ & $\sim$ & \texttimes & \texttimes \\
Arthur AI~\cite{arthurai2024} & $\sim$ & \texttimes & \texttimes & \texttimes & \texttimes & \texttimes & \texttimes & \texttimes & \texttimes \\
NIST Playbook~\cite{nistrmfplaybook} & \texttimes & \texttimes & $\sim$ & \texttimes & \texttimes & \texttimes & \texttimes & \texttimes & \texttimes \\
\midrule
\textbf{\regsynth{}} & \checkmark & \checkmark & \checkmark & \checkmark & $\sim$ & \checkmark & \checkmark & $\sim$ & \checkmark \\
\bottomrule
\end{tabular}
\end{table*}

% ============================================================================
\section{Evaluation}
\label{sec:eval}

Our current evaluation is intentionally limited to \emph{synthetic} compliance scenarios derived from public regulatory frameworks. The benchmark harness instantiates small collections of machine-encodable obligations, cost vectors, and phase-in schedules; it should be read as a prototype study rather than as a claim that the full legal content of those frameworks has been captured.

\subsection{Prototype Benchmark Scope}

The released scripts generate scenarios inspired by combinations of the EU AI Act, GDPR, CCPA, health-privacy rules, export controls, and related governance documents. These scenarios are useful because they permit controlled experiments over satisfiable and unsatisfiable formal fragments while keeping the optimization problems small enough to inspect by hand.

\subsection{Evidence from Reruns}

The current paper narrative is grounded only in the commands we reran against the repository during recovery.

\begin{itemize}[leftmargin=*]
  \item \texttt{cd implementation \&\& cargo test --release} rebuilt the workspace and executed the available release tests. The run exposed two remaining failures in the certificate crate (\texttt{pareto\_cert::tests::generate\_pareto\_cert\_two\_points} and \texttt{verifier::tests::verify\_pareto\_cert}), so the artifact currently supports a strong implementation claim but not a blanket claim that all certificate paths are validated end to end.
  \item \texttt{./implementation/target/release/regsynth benchmark --iterations 2 --num-obligations 10 --num-jurisdictions 2 --output-format json} completed successfully on a toy synthetic workload. Both iterations were feasible; the run reported an average of 14.5 generated constraints, frontier size 1.0, and stage timings rounded to 0\,ms.
  \item \texttt{python3 run\_benchmark.py} reran the bundled eight-scenario synthetic benchmark suite and refreshed \texttt{benchmarks/real\_benchmark\_results.json}. In this rerun, \texttt{MaxSMT\_Synthesis} succeeded on 5/8 scenarios (62.5\% success rate), achieved average hypervolume 0.3219, and preserved 100\% hard-constraint satisfaction on successful runs. However, \texttt{Epsilon\_Constraint} achieved the best average hypervolume in the same rerun (0.4828), followed by \texttt{Weighted\_Sum} (0.4068) and \texttt{NSGA-II} (0.3911). The MaxSMT path also failed on two scenarios because the current script still includes symbolic bounds that do not parse as floats.
  \item \texttt{python3 benchmarks/pareto\_dominance\_validation.py --runs 5 --bf-runs 10 --traj-runs 20} completed successfully and refreshed \texttt{benchmarks/results/pareto\_dominance\_validation.json}. The rerun reported 100\% dominance-correctness on the sampled checks, mean hypervolume ratio 0.9840 versus brute force (minimum 0.9158), and trajectory domination in 15\% of sampled instances with mean hypervolume improvement 21.84\%.
\end{itemize}

We inspected but did not use the historical high-dimensional benchmark outputs as grounding evidence in this paper revision. The default high-dimensional script is substantially more expensive to rerun, and its generated recommendation text overstates what the presently rerun evidence justifies.

\subsection{Prototype Findings}

These reruns support a narrower but still useful conclusion. \regsynth{} contains a real executable pipeline for synthetic multi-jurisdictional obligation sets, including encodings, trade-off exploration, and trajectory reasoning. The Pareto-validation harness provides the clearest executable support for the optimization story; the CLI benchmark demonstrates that the Rust toolchain is live; and the synthetic SOTA harness shows that the repository can compare optimization backends on bundled toy scenarios.

At the same time, the recovered evidence cuts against several stronger older claims. The current benchmark harness does \emph{not} show MaxSMT dominating alternative optimizers on average hypervolume, and the two scenario failures show that solver robustness is incomplete even within the synthetic suite. The temporal trajectory examples are best read as synthetic demonstrations explaining why myopic per-timestep optimization can miss better long-horizon trade-offs. Likewise, infeasibility detection is a core implemented capability of the tool, yet this version of the paper still does not report a stand-alone quantitative accuracy study for that component.

Overall, the evaluation supports a defensible message: \regsynth{} is an executable prototype for formalized fragments of regulatory reasoning, not yet a validated large-scale compliance engine. Several quantitative questions---especially solver robustness, scaling to much larger encodings, and legal-expert validation of the regulatory abstractions---remain future work.

\section{Related Work}
\label{sec:related}

\textbf{Deontic Logic and Normative Systems.}
Von Wright's foundational work on deontic logic~\cite{vonwright1951} established
the formal basis for reasoning about obligations and permissions.
Governatori~et~al.~\cite{governatori2013compliance} developed computational
compliance checking frameworks based on defeasible logic, handling exceptions and
rule conflicts within single jurisdictions. Sartor~\cite{sartor2005legal}
formalized legal reasoning with argumentation frameworks.
These works address single-jurisdiction formalization but do not tackle
multi-jurisdictional tractability, optimization, or temporal trajectory planning.

\textbf{Formal Methods for Regulatory Compliance.}
Hashmi~et~al.~\cite{hashmi2018compliance} surveyed formal approaches to legal
compliance checking, identifying a gap between formal methods capabilities
and practical regulatory complexity.
Breaux~and~Antón~\cite{breaux2008analyzing} extracted rights and obligations from
regulatory texts using natural language processing.
Maxwell~and~Antón~\cite{maxwell2012checking} developed methods for checking
regulatory compliance of requirements specifications.
Our work extends this line by addressing the multi-jurisdictional case with
optimization and temporal planning.

\textbf{Multi-Objective Optimization.}
Ehrgott~\cite{ehrgott2005} provides the foundational treatment of multicriteria
optimization. Papadimitriou and Yannakakis~\cite{papadimitriou2000} established
the computational complexity of multiobjective optimization.
Zitzler~et~al.~\cite{zitzler1999hypervolume} introduced the hypervolume indicator for
comparing Pareto frontier quality. Our contribution adapts these techniques to
the regulatory compliance domain via MaxSMT-based enumeration rather than
evolutionary algorithms.

\textbf{SMT and MaxSMT Solving.}
De~Moura and Bj{\o}rner~\cite{demoura2008z3} developed Z3, the most widely used
SMT solver. Barrett~et~al.~\cite{barrett2010smtlib} standardized the SMT-LIB
input format. Fu and Malik~\cite{fumalik2006} introduced iterative MaxSAT
algorithms. Liffiton and Sakallah~\cite{liffiton2008mus} developed efficient
MUS extraction algorithms. Our solver backend builds on these foundations
with regulatory-domain-specific optimizations.

\textbf{GRC Platforms and AI Governance Tools.}
The global GRC market is projected to reach \$64.6 billion by 2025~\cite{grcmarket2024}.
OneTrust~\cite{onetrust2024}, IBM OpenPages~\cite{ibmopenpages2024}, and
ServiceNow GRC~\cite{servicenow2024} lead the market but operate at the workflow
and checklist level without formal reasoning capabilities. TrustArc~\cite{trustarc2024}
provides strong privacy-regulation coverage (900+ templates for GDPR,
CCPA~\cite{ccpa2018}, LGPD~\cite{brazillgpd}) but does not address AI-specific
obligations. AI-specific governance platforms including
Credo AI~\cite{credoai2024}, Holistic AI~\cite{holisticai2024}, and
Arthur AI~\cite{arthurai2024} provide risk scoring, fairness auditing, and
compliance mapping but use heuristic rather than formal methods for conflict
detection and do not support multi-objective optimization or certificate
generation.

\textbf{AI Fairness Toolkits.}
IBM AI Fairness 360~\cite{ibmaif360} provides 70+ fairness metrics and 11 bias
mitigation algorithms. Microsoft's Responsible AI Toolbox~\cite{msrai2023}
integrates InterpretML, Fairlearn, and error analysis. Google Model
Cards~\cite{googlemodelcards} standardize model documentation. These toolkits
address important aspects of responsible AI (overlapping with EU AI Act Art.~10
data governance and Art.~13 transparency) but operate at the individual-model
level, do not encode regulatory obligations from multiple jurisdictions, and
cannot perform cross-framework conflict detection or compliance optimization.

\textbf{Regulatory Frameworks and Compliance Standards.}
The EU AI Act~\cite{euaiact2024} entered force in August 2024 as the first
comprehensive AI regulation, with phased enforcement through 2027. GDPR
Art.~22~\cite{gdpr2016} restricts automated decision-making for EU data subjects.
US Executive Order 14110~\cite{useo14110} establishes federal AI governance
requirements including safety testing and equity assessments. The
CCPA/CPRA~\cite{ccpa2018} provides California consumer privacy rights including
automated decision-making opt-out. Singapore's PDPA~\cite{singaporepdpa} and
Brazil's LGPD~\cite{brazillgpd} impose data protection obligations relevant to
AI systems. Canada's proposed AIDA (Bill C-27)~\cite{canadaaida} would create
binding AI governance requirements. The NIST AI RMF
Playbook~\cite{nistrmfplaybook} provides implementation guidance but is a
document-based framework rather than a computational tool.
Regulatory scholarship by Veale and Borgesius~\cite{veale2021demystifying}
analyzes the EU AI Act's implications for industry.
At the US state level, Colorado's SB~24-205~\cite{coloradoaiact2024} requires
algorithmic impact assessments for high-risk AI, California's SB~1047~\cite{californiasb1047}
mandates safety evaluations for frontier models, and Illinois
BIPA~\cite{illinoisbipa} restricts biometric data collection — creating a
patchwork of state-level requirements that compounds the multi-jurisdictional
challenge. EU AI Act compliance checkers~\cite{euaicheckers} have emerged
from law firms and compliance vendors but provide only decision-tree risk
classification without formal reasoning.

\textbf{Software Bill of Materials and Compliance Formats.}
The SPDX standard~\cite{spdx2023} provides a machine-readable format for
software package metadata and licensing information, increasingly relevant as
AI regulations require documentation of training data provenance and model
components. The OpenChain specification (ISO/IEC 5230:2020)~\cite{openchain2020}
establishes requirements for open source compliance programs, complementing
regulatory compliance with supply-chain transparency. \regsynth{} supports
export to both SPDX and JSON-LD formats for interoperability with existing
compliance toolchains.

\textbf{Semantic Web Technologies for Legal Knowledge.}
The European Legislation Identifier (ELI) provides a URI-based identification
system for European legislation, enabling linked-data representations of
regulatory provisions. OWL ontologies have been used to model legal concepts
and relationships, particularly in the Semantic Web community. \regsynth{}'s
export capabilities include OWL/RDF and JSON-LD formats, bridging the gap
between formal compliance analysis and the broader legal knowledge graph
ecosystem.

% ============================================================================
\section{Conclusion and Future Work}
\label{sec:conclusion}

We presented \regsynth{}, a prototype tool for formalizing a machine-readable fragment of multi-jurisdictional AI regulation and exploring trade-offs among candidate compliance strategies. Its main strengths today are architectural and executable rather than benchmark-superiority claims: a typed DSL, provenance-aware encodings, conflict diagnosis over the formalized fragment, multi-objective synthesis, temporal planning, and certificate generation inside a single Rust workspace.

The recovered evaluation supports a positive but narrower conclusion than earlier drafts. The repository builds, the CLI benchmark runs on small synthetic instances, the synthetic SOTA harness executes but shows that lighter baselines outperform the current MaxSMT path on average hypervolume, and the Pareto-validation harness supplies useful evidence for non-domination and trajectory-level reasoning. These results support \regsynth{} as an end-to-end prototype for synthetic regulatory fragments. They do \emph{not} support strong claims about infeasibility-detection accuracy, industrial-scale runtime, or superiority over alternative optimizers on large real regulatory workloads.

\textbf{Limitations.}
\regsynth{} operates on the formalizable fragment of regulatory obligations, which we estimate at 60--70\% of typical AI governance requirements. The remaining 30--40\% (graded as Opaque or F5) requires human judgment. Formalizability grading itself is subjective and has not been externally validated. The synthetic benchmark methodology is valuable for controlled experiments but does not substitute for legal-expert review or deployment-scale validation. In addition, the current release-test rerun still leaves two failing certificate tests, so certificate claims should be read as implementation scope rather than fully validated end-to-end behavior.

\textbf{Future work} includes: (1)~external validation of regulatory encodings by legal domain experts, (2)~improved solver robustness for richer encodings, including the scenario failures surfaced by the rerun benchmark harness, (3)~larger benchmark suites with measured scaling studies, (4)~sector-specific regulatory libraries, and (5)~incremental re-optimization when regulations change over time.

\appendix
\section{Proof of Theorem~\ref{thm:trajectory}: Trajectory Domination}
\label{app:proof}

We construct a concrete 3-timestep, 2-jurisdiction instance demonstrating that
per-timestep Pareto optimality does not imply trajectory-level Pareto optimality.

\textbf{Setup.} Consider two jurisdictions $J_1$ (EU) and $J_2$ (US), with
obligations activating at timesteps $t_1$, $t_2$, $t_3$. The transition budget
is $B = 1$ (at most one compliance action per timestep).

\textbf{Obligations.}
At $t_1$: Two obligations $O_1^{J_1}$ (cost $(1,3)$) and $O_1^{J_2}$ (cost $(3,1)$).
At $t_2$: Three obligations, where $O_2^{J_1}$ requires $O_1^{J_1}$ as prerequisite
(cost $(1,1)$ if $O_1^{J_1}$ active, $(4,4)$ otherwise).
At $t_3$: Final obligation $O_3$ with cost depending on trajectory history.

\textbf{Per-timestep optimal.} At $t_1$, both $(1,3)$ and $(3,1)$ are
Pareto-optimal; suppose the decision-maker selects the $(1,3)$-cost strategy
(choosing $O_1^{J_1}$). At $t_2$, with budget $B = 1$, the reachable strategies
from this state include $O_2^{J_1}$ at cost $(1,1)$. At $t_3$, the cumulative
trajectory cost is $(1+1+c_3, 3+1+c_3') = (2+c_3, 4+c_3')$.

\textbf{Jointly optimal.} Starting with $O_1^{J_2}$ at $t_1$ (cost $(3,1)$,
locally dominated on dimension 1), the trajectory through $t_2$ and $t_3$ achieves
cumulative cost $(3+1+c_3'', 1+1+c_3''')$ where $c_3'' < c_3$ and $c_3''' < c_3'$
due to the prerequisite structure. Specifically, choosing the ``worse'' option at
$t_1$ enables access to a cheaper compliance path at $t_2$ and $t_3$.

The full numerical construction yields $T^{\text{local}} = (4, 7)$ and
$T^{\text{global}} = (3, 5)$, demonstrating strict domination:
$T^{\text{global}} \prec T^{\text{local}}$. \qed

\section{DSL Grammar (Excerpt)}
\label{app:grammar}

\begin{lstlisting}[language=,caption={RegSynth DSL formal grammar (excerpt).}]
program     ::= declaration*
declaration ::= jurisdiction_decl | mapping_decl
                | import_decl

jurisdiction_decl ::=
  'jurisdiction' IDENT '{'
    framework_clause?
    obligation_decl*
  '}'

obligation_decl ::=
  'obligation' IDENT ':' modality
    '[' grade ']' '{'
    article_clause
    scope_clause?
    effective_clause?
    content_clause
    cost_clause
  '}'

modality    ::= 'OBL' | 'PERM' | 'PROH'
grade       ::= 'F1' | 'F2' | 'F3' | 'F4'
              | 'F5' | 'Opaque'
content     ::= expr
expr        ::= 'forall' IDENT 'in' path
                '=>' expr
              | 'exists' IDENT 'in' path
                '=>' expr
              | expr '&&' expr | expr '||' expr
              | 'not' expr  | comparison
              | IDENT '.' IDENT
comparison  ::= expr comp_op expr
comp_op     ::= '<=' | '>=' | '==' | '!='
                | '<' | '>'

cost_clause ::= 'cost' ':' '{'
  'financial' ':' NUMBER ','
  'time' ':' NUMBER ','
  'risk' ':' NUMBER ','
  'complexity' ':' NUMBER
'}'
\end{lstlisting}

% ============================================================================
% Bibliography
% ============================================================================

\begin{thebibliography}{40}

\bibitem{euaiact2024}
European Parliament and Council of the European Union.
\newblock Regulation (EU) 2024/1689 laying down harmonised rules on artificial
  intelligence (AI Act).
\newblock \emph{Official Journal of the European Union}, L series, 12 July 2024.

\bibitem{gdpr2016}
European Parliament and Council of the European Union.
\newblock Regulation (EU) 2016/679 on the protection of natural persons with
  regard to the processing of personal data (GDPR).
\newblock \emph{Official Journal of the European Union}, L 119, 4 May 2016.

\bibitem{nistrmf2023}
National Institute of Standards and Technology.
\newblock Artificial Intelligence Risk Management Framework (AI RMF 1.0).
\newblock NIST AI 100-1, January 2023.

\bibitem{iso42001}
International Organization for Standardization.
\newblock ISO/IEC 42001:2023 --- Information technology --- Artificial
  intelligence --- Management system.
\newblock 2023.

\bibitem{chinagai2023}
Cyberspace Administration of China.
\newblock Interim Measures for the Management of Generative Artificial
  Intelligence Services.
\newblock July 2023.

\bibitem{ieee7000}
IEEE.
\newblock IEEE 7000-2021 --- Standard Model Process for Addressing Ethical
  Concerns During System Design.
\newblock 2021.

\bibitem{vonwright1951}
G.~H. von Wright.
\newblock Deontic logic.
\newblock \emph{Mind}, 60(237):1--15, 1951.

\bibitem{governatori2013compliance}
G.~Governatori, F.~Olivieri, A.~Rotolo, and S.~Scannapieco.
\newblock Computing strong and weak permissions in defeasible logic.
\newblock \emph{Journal of Philosophical Logic}, 42(6):799--829, 2013.

\bibitem{sartor2005legal}
G.~Sartor.
\newblock \emph{Legal Reasoning: A Cognitive Approach to the Law}.
\newblock Springer, 2005.

\bibitem{hashmi2018compliance}
M.~Hashmi, G.~Governatori, H.~P. Lam, and M.~T. Wynn.
\newblock Are we done with business process compliance: state of the art and
  challenges ahead.
\newblock \emph{Knowledge and Information Systems}, 57(1):79--133, 2018.

\bibitem{breaux2008analyzing}
T.~D. Breaux and A.~I. Ant{\'o}n.
\newblock Analyzing regulatory rules for privacy and security requirements.
\newblock \emph{IEEE Transactions on Software Engineering}, 34(1):5--20, 2008.

\bibitem{maxwell2012checking}
J.~C. Maxwell and A.~I. Ant{\'o}n.
\newblock The production rule framework: developing a canonical set of software
  engineering production rules from regulatory texts.
\newblock In \emph{Proc. 20th IEEE International Requirements Engineering
  Conference (RE)}, pp. 51--60, 2012.

\bibitem{ehrgott2005}
M.~Ehrgott.
\newblock \emph{Multicriteria Optimization}.
\newblock Springer, 2nd edition, 2005.

\bibitem{papadimitriou2000}
C.~H. Papadimitriou and M.~Yannakakis.
\newblock On the approximability of trade-offs and optimal access of web
  sources.
\newblock In \emph{Proc. 41st Annual Symposium on Foundations of Computer
  Science (FOCS)}, pp. 86--92, 2000.

\bibitem{zitzler1999hypervolume}
E.~Zitzler and L.~Thiele.
\newblock Multiobjective evolutionary algorithms: a comparative case study and
  the strength Pareto approach.
\newblock \emph{IEEE Transactions on Evolutionary Computation}, 3(4):257--271,
  1999.

\bibitem{demoura2008z3}
L.~de~Moura and N.~Bj{\o}rner.
\newblock Z3: An efficient SMT solver.
\newblock In \emph{Proc. Tools and Algorithms for the Construction and Analysis
  of Systems (TACAS)}, LNCS 4963, pp. 337--340. Springer, 2008.

\bibitem{barrett2010smtlib}
C.~Barrett, A.~Stump, and C.~Tinelli.
\newblock The SMT-LIB standard: Version 2.0.
\newblock In \emph{Proc. 8th International Workshop on Satisfiability Modulo
  Theories (SMT)}, 2010.

\bibitem{fumalik2006}
Z.~Fu and S.~Malik.
\newblock On solving the partial MAX-SAT problem.
\newblock In \emph{Proc. International Conference on Theory and Applications of
  Satisfiability Testing (SAT)}, LNCS 4121, pp. 252--265. Springer, 2006.

\bibitem{liffiton2008mus}
M.~H. Liffiton and K.~A. Sakallah.
\newblock Algorithms for computing minimal unsatisfiable subsets of constraints.
\newblock \emph{Journal of Automated Reasoning}, 40(1):1--33, 2008.

\bibitem{necula1997pcc}
G.~C. Necula.
\newblock Proof-carrying code.
\newblock In \emph{Proc. 24th ACM Symposium on Principles of Programming
  Languages (POPL)}, pp. 106--119, 1997.

\bibitem{chaff2001}
M.~W. Moskewicz, C.~F. Madigan, Y.~Zhao, L.~Zhang, and S.~Malik.
\newblock Chaff: Engineering an efficient SAT solver.
\newblock In \emph{Proc. 38th Design Automation Conference (DAC)},
  pp. 530--535, 2001.

\bibitem{marques2009conflict}
J.~Marques-Silva, I.~Lynce, and S.~Malik.
\newblock Conflict-driven clause learning SAT solvers.
\newblock In \emph{Handbook of Satisfiability}, pp. 131--153. IOS Press, 2009.

\bibitem{onetrust2024}
OneTrust LLC.
\newblock OneTrust Privacy \& Data Governance Cloud.
\newblock \url{https://www.onetrust.com/}, accessed 2024.

\bibitem{ibmopenpages2024}
IBM Corporation.
\newblock IBM OpenPages with Watson.
\newblock \url{https://www.ibm.com/products/openpages-with-watson}, accessed 2024.

\bibitem{servicenow2024}
ServiceNow, Inc.
\newblock ServiceNow Governance, Risk, and Compliance (GRC).
\newblock \url{https://www.servicenow.com/products/governance-risk-and-compliance.html}, accessed 2024.

\bibitem{veale2021demystifying}
M.~Veale and F.~Zuiderveen~Borgesius.
\newblock Demystifying the Draft EU Artificial Intelligence Act.
\newblock \emph{Computer Law \& Security Review}, 41:105573, 2021.

\bibitem{grcmarket2024}
Grand View Research.
\newblock Governance, Risk and Compliance Market Size Report, 2024--2030.
\newblock Market Research Report GVR-4-68039-895-2, 2024.

\bibitem{achlioptas2005random}
D.~Achlioptas.
\newblock Random satisfiability.
\newblock In \emph{Handbook of Satisfiability}, pp. 245--270. IOS Press, 2009.

\bibitem{bellman1957}
R.~Bellman.
\newblock \emph{Dynamic Programming}.
\newblock Princeton University Press, 1957.

\bibitem{marco2016}
M.~H. Liffiton and A.~Malik.
\newblock Enumerating infeasibility: Finding multiple MUSes quickly.
\newblock In \emph{Proc. International Conference on Integration of Constraint
  Programming, Artificial Intelligence, and Operations Research (CPAIOR)},
  LNCS 7874, pp. 160--175, 2013.

\bibitem{coloradoaiact2024}
Colorado General Assembly.
\newblock SB 24-205: Concerning Consumer Protections for Artificial Intelligence.
\newblock Signed into law, May 2024.

\bibitem{californiasb1047}
California State Legislature.
\newblock SB-1047: Safe and Secure Innovation for Frontier Artificial
  Intelligence Models Act.
\newblock 2024.

\bibitem{illinoisbipa}
Illinois General Assembly.
\newblock 740 ILCS 14 --- Biometric Information Privacy Act (BIPA).
\newblock 2008.

\bibitem{spdx2023}
The Linux Foundation.
\newblock SPDX Specification v2.3.
\newblock \url{https://spdx.github.io/spdx-spec/v2.3/}, 2023.

\bibitem{openchain2020}
International Organization for Standardization.
\newblock ISO/IEC 5230:2020 --- OpenChain Specification.
\newblock 2020.

\bibitem{ibmaif360}
R.~K.~E.~Bellamy, K.~Dey, M.~Hind, S.~C.~Hoffman, S.~Houde,
  K.~Kannan, P.~Lohia, J.~Martino, S.~Mehta, A.~Mojsilovi{\'c}, et~al.
\newblock AI Fairness 360: An extensible toolkit for detecting and mitigating
  algorithmic bias.
\newblock \emph{IBM Journal of Research and Development}, 63(4/5):4:1--4:15,
  2019.

\bibitem{msrai2023}
Microsoft Corporation.
\newblock Responsible AI Toolbox: An open-source framework for building
  responsible AI solutions.
\newblock \url{https://responsibleaitoolbox.ai/}, accessed 2024.

\bibitem{googlemodelcards}
M.~Mitchell, S.~Wu, A.~Zaldivar, P.~Barnes, L.~Vasserman, B.~Hutchinson,
  E.~Spitzer, I.~D. Raji, and T.~Gebru.
\newblock Model cards for model reporting.
\newblock In \emph{Proc.\ ACM Conference on Fairness, Accountability, and
  Transparency (FAccT)}, pp. 220--229, 2019.

\bibitem{holisticai2024}
Holistic AI Ltd.
\newblock Holistic AI: AI Governance Platform.
\newblock \url{https://www.holisticai.com/}, accessed 2024.

\bibitem{credoai2024}
Credo AI Inc.
\newblock Credo AI: Responsible AI Governance Platform.
\newblock \url{https://www.credo.ai/}, accessed 2024.

\bibitem{arthurai2024}
Arthur AI Inc.
\newblock Arthur AI: AI Performance Monitoring.
\newblock \url{https://www.arthur.ai/}, accessed 2024.

\bibitem{trustarc2024}
TrustArc Inc.
\newblock TrustArc Privacy Management Platform.
\newblock \url{https://trustarc.com/}, accessed 2024.

\bibitem{onetrustai2024}
OneTrust LLC.
\newblock OneTrust AI Governance: Manage AI risk and compliance.
\newblock \url{https://www.onetrust.com/products/ai-governance/}, accessed 2024.

\bibitem{nistrmfplaybook}
National Institute of Standards and Technology.
\newblock NIST AI RMF Playbook.
\newblock \url{https://airc.nist.gov/AI_RMF_Playbook}, accessed 2024.

\bibitem{euaicheckers}
European Commission.
\newblock Assessment List for Trustworthy AI (ALTAI) --- Self-assessment tool
  for the EU AI Act.
\newblock \url{https://altai.insight-centre.org/}, accessed 2024.

\bibitem{useo14110}
The White House.
\newblock Executive Order 14110 on the Safe, Secure, and Trustworthy
  Development and Use of Artificial Intelligence.
\newblock Federal Register, 88 FR 75191, 30 October 2023.

\bibitem{ccpa2018}
California State Legislature.
\newblock California Consumer Privacy Act (CCPA), Cal.\ Civ.\ Code
  \S\S 1798.100--1798.199.100, as amended by CPRA (2020).
\newblock 2018.

\bibitem{singaporepdpa}
Personal Data Protection Commission, Singapore.
\newblock Personal Data Protection Act 2012 (PDPA), No.~26 of 2012, revised
  2021.
\newblock \url{https://www.pdpc.gov.sg/}, accessed 2024.

\bibitem{brazillgpd}
Federative Republic of Brazil.
\newblock Lei Geral de Prote{\c{c}}{\~a}o de Dados (LGPD), Law No.~13.709.
\newblock 14 August 2018, effective September 2020.

\bibitem{canadaaida}
Parliament of Canada.
\newblock Bill C-27: Digital Charter Implementation Act, 2022 --- Part 3:
  Artificial Intelligence and Data Act (AIDA).
\newblock First reading, 16 June 2022.

\end{thebibliography}

\end{document}
